\subsection{Resumen RAFA}

\begin{infobox}[cyan]{\textbf{Producción y detección de ondas gravitacionales primordiales durante el recalentamiento}}

El modelo cosmológico estándar incluye un período de \textit{inflación} que resuelve el problema del horizonte en el fondo cósmico de microondas (CMB). Tras esta etapa, el \textit{recalentamiento} conecta la inflación con la era de radiación, permitiendo posteriormente la nucleosíntesis primordial (BBN). Durante este período, el inflatón decae para generar la materia del universo mediante procesos que generan inhomogeneidades, los cuales a su vez producen ondas gravitacionales (GWs) \cite{figueroa_caprini, maggiore}.

En este trabajo estudiamos las características del espectro de GWs generado para distintos modelos de recalentamiento y los mecanismos físicos involucrados. Esto mismo, se realiza a partir de utilizar tanto estimaciones teóricas como los más actuales códigos de lattice para universos en expansión tal como es CosmoLattice \cite{cosmolattice} con el objetivo de:

\begin{enumerate}
    \item[(i)] Identificar los procesos físicos que generan perturbaciones tensoriales (GWs),
    
    \item[(ii)] Modelar su evolución tras entrar al horizonte cósmico, y
    
    \item[(iii)] Predecir el espectro de energía medido hoy $\left( \Omega_{\mathrm{gw}}(f) \right)$ y evaluar su viabilidad en detectores como LISA o redes de temporización de púlsares (PTAs).
\end{enumerate}

\begin{thebibliography}{9}

\bibitem{figueroa_caprini} 
C.~Caprini and D.~G.~Figueroa,
\textit{Cosmological Backgrounds of Gravitational Waves},
Class. Quantum Grav. \textbf{35} (2018) 163001,
\texttt{doi:10.1088/1361-6382/aac608},
arXiv:1801.04268 [astro-ph.CO].

\bibitem{maggiore} 
M.~Maggiore,
\textit{Gravitational Waves: Volume 2, Astrophysics and Cosmology},
Oxford University Press (2018).

\bibitem{cosmolattice} 
D.~G.~Figueroa, A.~Florio, F.~Torrenti, and W.~van~der~Schee, 
\textit{CosmoLattice: A modern code for lattice simulations of scalar and gauge field dynamics in an expanding universe}, 
Comput. Phys. Commun. \textbf{283} (2023) 108586, 
arXiv:2102.01031 [astro-ph.CO].

\end{thebibliography}

\end{infobox}